\documentclass[11pt,a4paper]{article}
\usepackage[margin=2.9cm]{geometry}
\usepackage{amsmath,amssymb,amsthm,mathrsfs}
\usepackage[T1]{fontenc}
\usepackage{microtype}
\usepackage{booktabs,longtable,array}
\usepackage[colorlinks=true,linkcolor=blue,citecolor=blue]{hyperref}

% ------------------------------------------------------------------
% Scoping pass for gate (II.11) of the Part II programme (strategy.md \S6-\S7).
% Statement inventory + design decisions, prior to execution. 2026-08-12.
% Proof-contract: every design decision is backed by a displayed computation;
% statements scheduled for execution are marked [exec]; imports are marked with
% candidate pinpoints. No scheduled statement is used as a premise elsewhere.
% ------------------------------------------------------------------

\newcolumntype{P}[1]{>{\raggedright\arraybackslash}p{#1}}

\theoremstyle{plain}
\newtheorem{theorem}{Theorem}[section]
\newtheorem{proposition}[theorem]{Proposition}
\newtheorem{lemma}[theorem]{Lemma}
\newtheorem{corollary}[theorem]{Corollary}
\theoremstyle{definition}
\newtheorem{definition}[theorem]{Definition}
\newtheorem{remark}[theorem]{Remark}
\newtheorem{statement}[theorem]{Statement}

\newcommand{\R}{\mathbb{R}}
\newcommand{\Z}{\mathbb{Z}}
\newcommand{\C}{\mathbb{C}}
\newcommand{\T}{\mathbb{T}}
\newcommand{\E}{\mathbb{E}}
\newcommand{\W}{\mathfrak{W}}
\newcommand{\HH}{\mathcal{H}}
\newcommand{\Om}{\Omega}
\newcommand{\eps}{\varepsilon}
\newcommand{\ip}[2]{\langle #1,#2\rangle}
\DeclareMathOperator{\supp}{supp}
\DeclareMathOperator{\dist}{dist}
\DeclareMathOperator{\spec}{spec}
\DeclareMathOperator{\Trace}{Tr}

\title{Gate (II.11): scoping pass\\[4pt]
\large Exact statement inventory for the interacting fixed-coarse limit,\\
the momentum-source obstruction, and the thermal two-torus route}
\author{}
\date{12 August 2026}

\begin{document}
\maketitle

\begin{abstract}
This document is the preparatory scoping pass for the single remaining gate of the
Part II programme: the interacting fixed-coarse characteristic-functional limit
(II.11). It records, with proofs where they are short and load-bearing:
(i) the \emph{minimal} logical requirements of the gate --- pointwise convergence in the
coarse symbol suffices; no uniformity over source sets is required by the proved Weyl
criterion; (ii) an exact obstruction which makes step (2) of the recorded attack plan
(\texttt{strategy.md} \S7) unimplementable as literally written: sharp-time momentum
sources are not Cameron--Martin admissible for the Euclidean measures on either side
(the relevant $H^{-1}$ integral diverges identically), so ``characteristic functionals
including the momentum component'' cannot be Gaussian source functionals; and (iii) a
repaired architecture --- the \emph{thermal two-torus route} --- in which both sides of
(II.11) are represented by traces over the compact space--time torus
$\T_t\times\T_L$, the momentum component enters through exact finite-mode Mehler
formulas and a smooth mean shift rather than through distributional sources, the
$t\to\infty$ and $N\to\infty$ limits are interchanged by a two-line spectral lemma, and
the previously feared $N$-uniform gap input is \emph{derived} (in eventual form) from
partition-function convergence instead of being assumed. The route is reduced to eight
statements A1--A8 with assigned proof methods; a fallback route B and its precise
defect are recorded; the import-acquisition list and the numeric falsification
checkpoint are specified. This document proves no part of (II.11); it fixes what
``proving (II.11)'' shall mean, statement by statement.
\end{abstract}

\tableofcontents

\section{The gate and its minimal logical requirements}\label{sec:gate}

Throughout, the model and notation are those of the Part II model sheet
(\texttt{lamport\_part\_ii\_}\allowbreak\texttt{model\_sheet.tex}): the spatial torus $\T_L=\R/(2L\Z)$;
dyadic lattices $\Lambda_N$ with spacing $\eps_N=2^{-N}\eps_0$; Weyl algebras
$\W_{N,L}$ over $\ell^2(\Lambda_N)$ with the symplectic form of the model sheet;
the $D4$ refinement embeddings $\alpha_M^N:\W_{M,L}\to\W_{N,L}$ and the continuum
embeddings $\beta_M$; the nearest-neighbour free Hamiltonian $H^{(0)}_{N,L}$ with
dispersion $\gamma_N(k)=(m^2+4\eps_N^{-2}\sin^2(\eps_Nk/2))^{1/2}$, $k\in\Gamma_N$; the
finite-torus Wick covariance $c_{N,L}$; the interacting operator
\[
 K_{N,L}:=H^{(0)}_{N,L}
 +\lambda\,\eps_N\!\!\sum_{x\in\Lambda_N}\!{:}\Phi_N(x)^4{:}_{c_{N,L}}
 \qquad(g_N\equiv1),
\]
realized by its coercive closed quadratic form, with compact resolvent and simple
ground state $\Om_{N}$, and the states
$\omega_{N}:=\omega_{N,L,\lambda,1}=\ip{\Om_N}{\cdot\,\Om_N}$.

\begin{definition}[the gate]\label{def:gate}
Gate (II.11) is the statement: there is a state $\omega^{\rm ct}_{L,\lambda}$ on the
continuum Weyl algebra $\W(\T_L)$ --- namely the vector state of the ground vector
of the operator $K_L$ constructed in V3, Theorem 4.1 as the strong limit of the
mode-truncated semigroups, in the convention $(a,b)=(0,0)$, the plain $\varphi^4$
truncation convention (V2, Corollary 5.2; V3, \S1) --- such that for every
fixed coarse scale $M$ and every $\xi\in h_{M,L}$,
\begin{equation}
 \lim_{N\to\infty}\omega_{N}\bigl(\alpha_M^N(W_M(\xi))\bigr)
 =\omega^{\rm ct}_{L,\lambda}\bigl(\beta_M(W_M(\xi))\bigr).
 \tag{II.11}\label{eq:II11}
\end{equation}
Whether this $K_L$ coincides with any \emph{other} realization of
$H_{0,L}+\lambda\int_{\T_L}{:}\varphi^4{:}$ on $\T_L$ --- in particular with a Segal
essentially-self-adjoint one --- is neither proved nor consumed anywhere in Part~II;
see V3, Remark 4.2, which records that essential self-adjointness of
$(H_0+V_L)|_{\mathcal D}$ is not established here and that $\omega^{\rm ct}_{L,\lambda}$ is
\emph{defined} by this $K_L$. The gate is therefore an identification with a
constructed object, not a uniqueness claim.
\end{definition}

\begin{remark}[what the Weyl criterion already delivers; demotion of uniformity]
\label{rem:pointwise}
The full-sequence Weyl criterion and the identification corollary are \emph{proved}
(\texttt{lamport\_part\_ii\_}\allowbreak\texttt{weyl\_reduction.tex},
Thm.~1.1 and Cor.~3.1): if the limit
in \eqref{eq:II11} exists for every fixed $M$ and $\xi$, then the limits define states
$\rho_M$ on each $\W_{M,L}$, convergence upgrades automatically from Weyl generators to
every $A\in\W_{M,L}$ (a $2\delta$ argument on the norm-dense span), the family is
projectively consistent, and it assembles to the unique inductive-limit state; if
moreover the limit equals the right side of \eqref{eq:II11} generator by generator,
the inductive-limit state is $\omega^{\rm ct}_{L,\lambda}\circ\beta$. The hypotheses
are \emph{pointwise in $\xi$}. Consequently:
\begin{enumerate}
\item Step (3) of the recorded attack plan (\texttt{strategy.md} \S7: ``prove
convergence of those scalar functionals uniformly on a stated bounded set of source
strengths'') demands more than the gate needs. Uniformity on bounded symbol sets is a
\emph{desirable output} of the estimates below (and would feed any later quantitative
rate), but it is not a logical prerequisite and shall not be treated as one.
\item No positivity, normalization, or consistency statement needs separate proof at
the gate: all are downstream of pointwise convergence, by the cited criterion.
\end{enumerate}
\end{remark}

\begin{remark}[what is off the critical path --- reaffirmed with one refinement]
\label{rem:offpath}
The corrected programme records that an $N$-uniform interacting gap, $N$-uniform local
moments, and an interacting Lieb--Robinson bound are not on the static critical path.
The scoping pass refines this: the route below \emph{derives} an eventual uniform gap
(Statement~A7, via Lemma~\ref{lem:laplace}) from partition-function convergence. It is
an output, not an input --- the same logical pattern as the Part I slab limit, where
the gap of the cutoff theory was an output of CE2 and simplicity was the only
spectral input.
\end{remark}

\section{The momentum-source obstruction and the design amendment}
\label{sec:obstruction}

Step (2) of the recorded attack plan asks to ``represent the two sides of (II.11) by
their corresponding finite-torus Euclidean characteristic functionals, including the
momentum component of $\xi$.'' The following computation shows that the momentum
component admits \emph{no} representation as a Gaussian source functional, on either
side, so step (2) is unimplementable as written and must be amended.

\begin{lemma}[sharp-time momentum sources are not $H^{-1}$]\label{lem:nogo}
Let $C_\infty=(-\partial_t^2-\partial_x^2+m^2)^{-1}$ be the free Euclidean covariance
on $\R_t\times\T_L$, and for real $h\in C^\infty(\T_L)$, $h\ne0$, let
$\delta_0'\otimes h$ be the time-derivative sharp-time source, i.e.\ the functional
$\varphi\mapsto-\partial_t\Phi(\delta_0\otimes h)$ in the distributional sense. Then
\begin{equation}
 \bigl\|\delta_0'\otimes h\bigr\|_{H^{-1}}^2
 =\frac1{2L}\sum_{k\in\Gamma_\infty}|\widehat h(k)|^2
 \int_\R\frac{dk_0}{2\pi}\,\frac{k_0^2}{k_0^2+\gamma(k)^2}
 =+\infty ,
 \label{eq:nogo}
\end{equation}
where $\gamma(k)^2=k^2+m^2$. The same divergence, with $\gamma_N$ in place of
$\gamma$, occurs for the Euclidean measure of the spatial-lattice model (whose time
variable is continuous). Consequently the random variable ``$\Phi(\delta_0'\otimes h)$''
is not an element of the Gaussian first chaos on either side, and
$e^{i\Pi(h)}$-insertions cannot be written as $e^{i\Phi(J)}$ with an $H^{-1}$ source
$J$.
\end{lemma}

\begin{proof}
The Fourier symbol of $\delta_0'\otimes h$ is $ik_0\widehat h(k)$, so the
$H^{-1}$-norm is the displayed double sum; the $k_0$-integrand tends to $1$ as
$|k_0|\to\infty$ for every fixed $k$, so already a single $k$-term with
$\widehat h(k)\neq0$ diverges. For the lattice model the spatial sum is finite but the
$k_0$-integral is unchanged.
\end{proof}

\begin{remark}[design amendment D-$\alpha$]\label{rem:amendment}
The momentum component of $\xi=\eps_Mq+ip$ shall never enter path space as a source.
It enters in exactly two controlled ways:
\begin{enumerate}
\item \emph{Free/thermal level (exact formulas).} All free thermal Weyl expectations
are computed in closed form by finite-mode Mehler algebra (Statement~A3); no path-space
representation is used for them.
\item \emph{Interacting level (mean shift).} In the trace representation of
Statement~A2, the unitary $e^{i\Pi_N(p)}$ acts by translation of the field
configuration at the marked time slice. The resulting ``jump'' ensemble is the
translate of the periodic Gaussian ensemble by the \emph{explicit mean path}
$H_p(t')$ --- the free covariance's harmonic interpolation of the jump around the time
circle --- which is smooth away from the marked slice and lies in every
$L^p$ of the circle. The interaction then changes by the Wick-binomial finite sum
\begin{equation}
 \mathcal U(\varphi+H_p)-\mathcal U(\varphi)
 =\lambda\!\int\!\Bigl[\,4H_p\,{:}\varphi^3{:}+6H_p^2\,{:}\varphi^2{:}
 +4H_p^3\,\varphi+H_p^4\Bigr],
 \label{eq:meanshift}
\end{equation}
a polynomial perturbation with smooth deterministic coefficients, uniformly
controlled in $N$ because $H^{(N)}_p\to H_p$ by the dispersion estimates of the
free-rate document. No distributional source ever appears. The precise bookkeeping is
Statement~A2's content and is scheduled, not assumed.
\end{enumerate}
This amendment replaces step (2) of \texttt{strategy.md} \S7 and is recorded there.
\end{remark}

\section{Route A: the thermal two-torus architecture}\label{sec:routeA}

\subsection{The representation and the two elementary lemmas}

For $t>0$ define the thermal functionals
\begin{equation}
 F_N(t;\xi):=\frac{\Trace\bigl(e^{-tK_{N,L}}\,W_N(R_M^N\xi)\bigr)}
                  {\Trace\bigl(e^{-tK_{N,L}}\bigr)},
 \qquad
 F_\infty(t;\xi):=\frac{\Trace\bigl(e^{-tK_{L}}\,W_{\rm ct}(R_M^\infty\xi)\bigr)}
                  {\Trace\bigl(e^{-tK_{L}}\bigr)},
 \label{eq:thermal}
\end{equation}
the right-hand version conditional on the continuum package A1 (trace-class
semigroup). Both numerators are Euclidean objects over the \emph{compact} space--time
torus $\T_t\times\T_L$ (time circle of circumference $t$): no slab limit, no boundary
terms, no spatial cutoff. The architecture rests on the following two lemmas, which we
prove now because the entire route stands on them.

\begin{lemma}[interchange of $t\to\infty$ and $N\to\infty$]\label{lem:interchange}
Let $(a_N(t))_{N\in\mathbb N\cup\{\infty\}}$, $t\ge t_0>0$, be complex functions with
$|a_N(t)|\le c_0$ for all $N,t$, and let $\ell_N$ be numbers such that
\begin{equation}
 |a_N(t)-\ell_N|\le c_0\,D_N(t),\qquad
 D_N(t):=\sum_{j\ge1}e^{-t\,(E_j(N)-E_0(N))},
 \label{eq:tail}
\end{equation}
for discrete sequences $E_0(N)<E_1(N)\le E_2(N)\le\cdots\to\infty$. Assume
\begin{enumerate}
\item[\rm(i)] $a_N(t)\to a_\infty(t)$ as $N\to\infty$, for every fixed $t\ge t_0$;
\item[\rm(ii)] $D_N(t_1)\to D_\infty(t_1)<\infty$ as $N\to\infty$, for one fixed
$t_1\ge t_0$;
\item[\rm(iii)] $\liminf_N\,(E_1(N)-E_0(N))\ge\gamma>0$.
\end{enumerate}
Then $\ell_N\to\ell_\infty$.
\end{lemma}

\begin{proof}
For $t\ge t_1$,
\[
 D_N(t)=\sum_{j\ge1}e^{-(t-t_1)(E_j(N)-E_0(N))}\,e^{-t_1(E_j(N)-E_0(N))}
 \le e^{-(t-t_1)(E_1(N)-E_0(N))}\,D_N(t_1).
\]
By (ii) there are $N_0$ and $D_*<\infty$ with $D_N(t_1)\le D_*$ for
$N_0\le N<\infty$ and also $D_\infty(t_1)\le D_*$, and by (iii) there is
$N_1\ge N_0$ with $E_1(N)-E_0(N)\ge\gamma/2$ for $N_1\le N<\infty$. Hence
\[
 \sup_{N_1\le N<\infty}D_N(t)\le D_*\,e^{-(t-t_1)\gamma/2}
 \xrightarrow[t\to\infty]{}0 .
\]
Hypothesis (iii) is a statement about finite $N$ only and says nothing about
the limit member, which is handled separately by the standing strict
inequality $E_0(\infty)<E_1(\infty)$ of the setup:
$D_\infty(t)\le e^{-(t-t_1)(E_1(\infty)-E_0(\infty))}D_\infty(t_1)\to0$ as
$t\to\infty$ as well.
Given $\eps>0$ choose $t$ with $\sup_{N_1\le N<\infty}D_N(t)<\eps$ and
$D_\infty(t)<\eps$; then for $N\ge N_1$,
\[
 |\ell_N-\ell_\infty|
 \le|\ell_N-a_N(t)|+|a_N(t)-a_\infty(t)|+|a_\infty(t)-\ell_\infty|
 \le 2c_0\eps+|a_N(t)-a_\infty(t)| ,
\]
and the middle term tends to $0$ by (i). Thus
$\limsup_N|\ell_N-\ell_\infty|\le2c_0\eps$ for every $\eps$.
\end{proof}

\begin{lemma}[Laplace transfer of spectra]\label{lem:laplace}
For $N\in\mathbb N\cup\{\infty\}$ let $\nu_N=\sum_j\delta_{E_j(N)}$ be the eigenvalue
point measures (with multiplicity) of self-adjoint operators with
$Z_N(t):=\int e^{-tE}\,d\nu_N(E)<\infty$ for all $t>0$. Assume
\begin{enumerate}
\item[\rm(i)] $Z_N(t)\to Z_\infty(t)$ as $N\to\infty$, for every $t\ge t_*$, where
$t_*>0$ is fixed but arbitrary;
\item[\rm(ii)] $\inf_NE_0(N)>-\infty$.
\end{enumerate}
Then $\nu_N\to\nu_\infty$ vaguely on $\R$; consequently, for every $j$,
$E_j(N)\to E_j(\infty)$ (eigenvalues in increasing order with multiplicity), and if
$E_0(\infty)$ is simple with gap $\gamma_L:=E_1(\infty)-E_0(\infty)>0$, then
\[
 E_1(N)-E_0(N)\longrightarrow\gamma_L,
 \qquad
 D_N(t)=e^{tE_0(N)}Z_N(t)-1\longrightarrow D_\infty(t)\quad(t\ge t_*).
\]
\end{lemma}

\begin{proof}
Let $c:=\inf_NE_0(N)$. All $\nu_N$ are supported in $[c,\infty)$. The linear span
$\mathscr A_{t_*}$ of $\{E\mapsto e^{-tE}:t\ge t_*\}$ is a subalgebra of
$C_0([c,\infty))$ (indeed $e^{-tE}e^{-sE}=e^{-(t+s)E}$ with $t+s\ge2t_*\ge t_*$)
which separates points and vanishes nowhere, so by the Stone--Weierstrass theorem
for locally compact spaces $\mathscr A_{t_*}$ is uniformly dense in
$C_0([c,\infty))$. For $f\in\mathscr A_{t_*}$, (i) gives
$\int f\,d\nu_N\to\int f\,d\nu_\infty$.

The passage to general $f$ must be made against a \emph{finite} measure: $\nu_N$
has infinite total mass, so $\|f-g\|_\infty<\eps$ by itself controls nothing.
(For $E_j(N)=j$, $f(E)=\eps/(1+E)\in C_0$ and $g=0$ one has
$\int f\,d\nu_N=+\infty$.) We therefore pair against $e^{-t_*E}d\nu_N$, whose total
mass is $Z_N(t_*)$ and hence uniformly bounded by (i). Let $f\in C_c(\R)$ and
$\eps>0$. Then $\widetilde f(E):=f(E)e^{t_*E}\in C_c([c,\infty))\subset
C_0([c,\infty))$; pick $g\in\mathscr A_{t_*}$ with
$\|\widetilde f-g\|_\infty<\eps$. Since $\int e^{-t_*E}d\nu_N=Z_N(t_*)$,
\[
 \Bigl|\int f\,d\nu_N-\int g\,e^{-t_*E}d\nu_N\Bigr|
 =\Bigl|\int(\widetilde f-g)\,e^{-t_*E}d\nu_N\Bigr|\ \le\ \eps\,Z_N(t_*),
\]
and likewise for $\nu_\infty$. Moreover $g\,e^{-t_*E}\in\mathscr A_{t_*}$, so
$\int g\,e^{-t_*E}d\nu_N\to\int g\,e^{-t_*E}d\nu_\infty$ by (i). Letting
$\eps\downarrow0$ gives $\int f\,d\nu_N\to\int f\,d\nu_\infty$ for every
$f\in C_c(\R)$, which is vague convergence.

The restriction $t\ge t_*$ in the last conclusion is necessary, not cosmetic:
hypothesis (i) says nothing about $Z_N$ on $(0,t_*)$. Take $t_*=1$,
$\nu_\infty=\sum_{j\ge0}\delta_j$ and $\nu_N=\nu_\infty+m_N\delta_N$ with
$m_N:=\lceil e^N/N\rceil$. Then $Z_N(t)<\infty$ for every $t>0$; for $t\ge1$ one has
$m_Ne^{-tN}\le(e^N/N+1)e^{-N}\to0$, so (i) holds on $[t_*,\infty)$; $E_0(N)=0$, so (ii)
holds; and $E_j(N)=j\to E_j(\infty)$, so the eigenvalue and gap conclusions hold. Yet
$D_N(t)-D_\infty(t)=m_Ne^{-tN}$, which at $t=\tfrac12$ is $\sim e^{N/2}/N\to\infty$.
Only $t\ge t_*$ is claimed, and only one value $t_1\ge t_*$ is consumed downstream
(V4, Theorem 5.2(c)).

Vague convergence of integer-valued atomic measures to an atomic limit gives
convergence of the ordered atoms with multiplicity on every compact interval whose
endpoints are not atoms of $\nu_\infty$; since $\nu_\infty$ is locally finite and its
atoms are discrete, $E_j(N)\to E_j(\infty)$ for every $j$. The gap statement is the
case $j=0,1$; the convergence of
$D_N(t)=e^{tE_0(N)}Z_N(t)-1$ follows from $E_0(N)\to E_0(\infty)$ and (i).
\end{proof}

\begin{remark}
Lemma~\ref{lem:interchange} with $a_N=F_N(\cdot\,;\xi)$, $\ell_N=\omega_N(W)$,
$c_0=2$, and the spectral tail bound
$|F_N(t;\xi)-\omega_N(W_N(R_M^N\xi))|
=\bigl|Z_N(t)^{-1}\sum_{j\ge1}e^{-tE_j(N)}(\ip{\phi_j}{W\phi_j}-\omega_N(W))\bigr|
\le2D_N(t)/(1+D_N(t))\le2D_N(t)$
reduces \eqref{eq:II11} to: fixed-$t$ convergence of the thermal functionals
(A6), fixed-$t$ convergence of partition functions feeding
Lemma~\ref{lem:laplace} (A7), and the continuum package (A1). This is the entire
logical skeleton of Route A; everything else below is the analytic content of those
three items.
\end{remark}

\subsection{Statement inventory A1--A8}

\begin{statement}[A1: continuum torus package]\label{st:A1}
\emph{[import + exec]} For $\lambda\ge0$ and the convention of A4b: the continuum
operator $K_L=H_{0,L}+\lambda\int_{\T_L}{:}\varphi^4{:}_{C_L}\,dx$ is self-adjoint and
bounded below (form sum); $e^{-tK_L}$ is trace class for every $t>0$ (hence discrete
spectrum, $Z_\infty(t)<\infty$); the ground state is simple with gap
$\gamma_L>0$; and $\omega^{\rm ct}_{L,\lambda}$ is its vector state on
$\W(\T_L)$. Trace-class of the free part is elementary
($\Trace e^{-tH_{0,L}}=\prod_{k\in\Gamma_\infty}(1-e^{-t\gamma(k)})^{-1}<\infty$
since $\sum_ke^{-t\gamma(k)}<\infty$); the interacting statements are classical
finite-volume $P(\varphi)_2$ material. \emph{Import candidates:}
H{\o}egh-Krohn's finite-temperature $P(\varphi)_2$ paper (CMP 38 (1974) 195--224) and
Figari--H{\o}egh-Krohn--Nappi (CMP 44 (1975) 265--278, $P(\varphi)_2$ on the de Sitter
cylinder $=$ spatial circle), plus the $Q$-space semiboundedness/hypercontractivity
package of Simon's book as an import candidate. \emph{Fallback:} prove from scratch on
the torus by the Part I Appendix-C pattern (Nelson bounds on compact space are
strictly easier than on $\R$); simplicity by positivity improvement; trace-class by
Golden--Thompson against $e^{-tH_{0,L}}$ plus $e^{-\mathcal U}\in L^p$.
\emph{Deliverable:} a source-scope audit in the Part II dossier style deciding
import vs.\ proof, with page images for whatever is imported.
\emph{Executed resolution (16 August 2026):} V3 takes the proof route.  It derives
the finite-mode Gaussian unitary, closed form, compact embedding, Mehler--Trotter
kernel, cyclic trace, mean shift, periodic cutoff tails, and positivity argument
explicitly.  Its only printed import is the precisely enumerated $\Gamma(A)$ cluster
from Simon, with the consumed pages imaged in the dossier; no unproved transplantation
from a line model to the circle is used.
\end{statement}

\begin{statement}[A2: trace representations with a marked slice]\label{st:A2}
\emph{[exec]} (a) Lattice: for the finite-dimensional Schr\"odinger operator
$K_{N,L}$ and $W_N(\eps_Nq+ip)=e^{i(\Phi_N(q)+\Pi_N(p))}$, the kernel identity
\[
 \Trace\bigl(e^{-tK_{N,L}}W_N(\eps_Nq+ip)\bigr)
 =e^{i\theta(q,p)}\!\!\int_{\R^{\Lambda_N}}\!\!
 e^{-tK_{N,L}}\bigl(\varphi+p,\varphi\bigr)\,
 e^{i\,\eps_N\sum_xq(x)\varphi(x)}\,d\varphi ,
\]
with the explicit CCR phase $\theta$, and its Feynman--Kac form: the periodic
Ornstein--Uhlenbeck ensemble on $\T_t\times\Lambda_N$ with a configuration jump at the
marked slice, equal to the translate of the periodic ensemble by the explicit mean
path $H^{(N)}_p$, with the interaction shifted per \eqref{eq:meanshift}. All objects
are finite-dimensional; the proof is Trotter plus Gaussian bridge algebra.
(b) Continuum: the same statement over $\T_t\times\T_L$, with the mean path
$H_p(t')$ given by the periodized kernel
$\cosh\bigl((\tfrac t2-|t'|)\gamma\bigr)/\!\sinh(\tfrac t2\gamma)$-type mode
multipliers applied to $p$; proof by A1's trace-class semigroup, cylinder
approximation, and the same bridge algebra mode by mode. \emph{Contract:} part (b) is
formulated so that its free case is Statement A3's formula, and no object beyond A1 is
consumed.
\end{statement}

\begin{statement}[A3: exact free thermal functionals and a thermal fixed-coarse
rate]\label{st:A3}
\emph{[exec]} (a) Closed forms: for the free lattice and continuum models,
\[
\begin{gathered}
 F^{(0)}_N(t;\xi)=\exp\Bigl[-\tfrac14\,Q^{(t)}_{N}(\xi)\Bigr],\\
 Q^{(t)}_{N}(\xi)=\frac1{2r_N}\sum_{k\in\Gamma_N}
 \coth\!\Bigl(\tfrac{t\,\gamma_N(k)}2\Bigr)
 \Bigl(\gamma_N(k)^{-1}|\widehat q(k)|^2+\gamma_N(k)|\widehat p(k)|^2\Bigr),
\end{gathered}
\]
and the analogous $Q^{(t)}_{\infty}$ with $\gamma$, obtained per mode from the
harmonic-oscillator identity
$\langle W(z)\rangle_{t}=\exp[-\tfrac12|z|^2\coth(t\gamma/2)]$ (finite-mode Mehler
computation; to be derived, not imported). At $t\to\infty$, $\coth\to1$ recovers the
vacuum formulas of the free-rate document, eq.~(1.4)/(1.8).
(b) Thermal fixed-coarse rate: for fixed $M$ and $\xi$, uniformly in $t\ge t_0$,
\[
 \bigl|Q^{(t)}_{M+r}(R_M^{M+r}\xi)-Q^{(t)}_{\infty}(R_M^\infty\xi)\bigr|
 \le C\,2^{-\theta r},
\]
by the same envelope/filter-tail/dispersion decomposition as the vacuum rate theorem,
using $1\le\coth(t\gamma/2)\le\coth(t_0m/2)$ and, for $t\ge t_0$, $\gamma\ge m$,
\[
 \coth(t\gamma/2)-1=\frac2{e^{t\gamma}-1}\ \le\ \frac2{e^{t_0\gamma}-1}
 \ \le\ \frac{2e^{-t_0\gamma}}{1-e^{-t_0m}}
\]
(the naive bound $2e^{-t_0\gamma}$ is \emph{false} for every $t_0,\gamma>0$, since
$2/(e^x-1)-2e^{-x}=2/[e^x(e^x-1)]>0$, and the ratio blows up like $1/(t_0\gamma)$ as
$t_0\gamma\downarrow0$; only the corrected form above is available). The thermal
weights improve, never damage, the high-momentum tails. \emph{Contract:} (b) reuses the proved Lemmas 2.1--3.1 of the
free-rate document verbatim; only the weight bookkeeping is new.
\end{statement}

\begin{statement}[A4: interaction comparison on the space--time torus]\label{st:A4}
\emph{[executed form, repaired 16 August 2026]} Use one Hermitian Gaussian Fourier
family on $D_\infty$.  Off the Nyquist column set $G_N=G$; at
$k_e=-\pi/\eps_N$ set
$G_N(k_0,k_e)=2^{-1/2}[G(k_0,k_e)+G(k_0,-k_e)]$.  This is the V2
same-noise coupling and makes the lattice field real.  Its mixed covariance has a
translation-invariant part with symbol
$\hat c_\times=(\hat c_N\hat c_\infty)^{1/2}$ and a separately bounded physical-edge
remainder. Then:
\begin{enumerate}
\item[\rm(a)] \emph{$L^2$ convergence:}
$\|\mathcal U_N-\mathcal U\|_{L^2(\mathbb P)}
\le C(1+t)(1+L)\eps_N[1+\log(1/\eps_N)]^{3/2}$ for $tL\ge1$.
The executed fourth-chaos identity is
\[
 \|X_4-Y_4\|_2^2=4!(2Lt)^{-2}B_4+\mathrm{Al}_N
 -2\mathcal E_N^{\rm edge}.
\]
Here telescoping the principal products gives the two explicit sums
$T_1\le C(tL)^3\eps_N^2\log_N^3$ (dispersion) and
$T_2\le C(tL)^3\eps_N^2\log_N^2$ (missing inputs), while the three possible
lattice alias outputs satisfy
$\mathrm{Al}_N\le C tL\eps_N^2\log_N^3$.  These are the bounds proved from
V2's displayed $F_3$ convolution estimate and one-dimensional temporal and
spatial sum--integral comparisons; no power-counting slogan is used as a proof.
At the symmetric Nyquist edges,
four momenta can conserve exactly (two positive and two negative); their missing
contribution is
\[
 \Delta_{4e}=4!(2Lt)^{-2}\frac32
 \sum_{\sum k_{0,i}=0}\prod_i\hat c_\times(k_{0,i},\pi/\eps_N)
 \le\frac{9t}{256L^2}\coth(2t_0)^3\eps_N^5.
\]
This term is retained in the executed V2 proof.
\item[\rm(b)] \emph{Finite renormalization dictionary (A4b):} the model sheet fixes
Wick ordering by the \emph{lattice vacuum} constant $c_{N,L}$, while the continuum
interaction is Wick-ordered by a continuum convention $c^{\rm ct}_{L}$ (to be fixed:
vacuum mode sum with a stated truncation). The recombination identity
${:}\varphi^4{:}_{c_1}={:}\varphi^4{:}_{c_2}+6(c_2-c_1){:}\varphi^2{:}_{c_2}
+3(c_2-c_1)^2$ shows: if
$\delta_N:=c_{N,L}-c^{\rm ct}_{L,N}$ (matched truncations) converges to
$\delta_\infty$, then --- \emph{if the two Wick powers were taken of the same field}
--- the lattice theory would converge to the continuum theory with polynomial
$\varphi^4-6\delta_\infty\varphi^2+3\delta_\infty^2$ (note the sign, which follows from
$c_2-c_1=-\delta_N$ in the displayed identity; the version with $+6\delta_\infty$
printed here originally was a slip).
\emph{Superseded (12 Aug 2026):} the executed form of this clause is V2,
Corollary 5.2, and its outcome is that \emph{no} residue survives: the lattice and
continuum fields' equal-point variances also differ by exactly $\delta_\infty$, which
cancels the excess Wick subtraction, so the gate's continuum polynomial is plain
$\varphi^4$ and the endpoint convention is $(a,b)=(0,0)$. $\delta_\infty$ remains the
correctly computed offset between the two tadpoles. \emph{Task (as executed):}
compute $\delta_\infty$ exactly (mode sums; the model sheet's
$c_N=\tfrac1{2\pi}\log\tfrac8{m\eps_N}+o(1)$ already fixes the divergent part), and
\emph{define} $\omega^{\rm ct}_{L,\lambda}$ in Definition~\ref{def:gate} as the vacuum
of the so-normalized continuum polynomial. This is a convention lock, exactly parallel
to Part I's Remark on sharp-time Wick ordering; (II.11) without it is
under-specified.
\item[\rm(c)] \emph{Mean-shift terms:} the same $L^2$ statements for each term of
\eqref{eq:meanshift}, with the deterministic coefficients $H^{(N)}_p\to H_p$ supplied
by A2/A3 dispersion estimates.  The executed proof uses the fundamental band
$B_N=\Gamma_N$ and representative $[q]_N$ modulo $2\pi/\eps_N$.  For chaos order
$j=1,2,3$ it distinguishes the folded lattice multiplier
$K_{N,j}^{\rm fold}$ from the exact-output multipliers $K_{\times,j}$ and
$K_{\infty,j}$.  Because the asymmetric-band mixed kernel need not be real on
the Nyquist edge, V2 defines its scalar surrogate as the real part of the mixed
pairing.  After extending the lattice coefficient by zero, that real surrogate
quadratic difference is exactly
\[
 \langle\delta b,K_{\infty,j}\delta b\rangle
 +\langle b_N,(K_{N,j}^{\rm fold}-K_{\infty,j})b_N\rangle
 -2\Re\langle b_N,(K_{\times,j}-K_{\infty,j})b_{\infty,B_N}\rangle
 -2\Re A_{\times,j},
\]
where $A_{\times,j}$ is the explicitly summed set of nonzero exact outputs
$q=p+2n\pi/\eps_N$.  Fundamental-band differences are
$O(\eps_N^2\log_N^2)$, high outputs are
$O(\eps_N^2\log_N^{j-1})$, and coefficient powers are controlled by a written
convolution telescope.  The physical edge remainder is estimated separately.
There is no unnamed tail term.  The resulting rate is
$O(2^{-r(1+\eta)/4}+\eps_N\log_N)$.
\end{enumerate}

\emph{Superseded (13 Aug 2026):} the coupling scheduled here --- the image of the
continuum space--time field under a sharp-time per-spatial-mode amplitude ratio --- is
\emph{not} the coupling executed, and would not do what it is asked to do: a per-mode
multiplier applied at a single time does not reproduce the lattice ensemble's time
correlations. What is executed is V2's \emph{same-noise} coupling (V2,
Definition 1.1): one Hermitian Gaussian family $G$ on $D_\infty$, with $\eta$ built
from the symbol $\hat c_\infty^{1/2}$ over $D_\infty$ and $\eta_N$ from
$\hat c_N^{1/2}$ over $D_N$, the edge column $k_e=-\pi/\eps_N$ symmetrised as
$G_N(k_0,k_e)=2^{-1/2}(G(k_0,k_e)+G(k_0,-k_e))$ (without which $\eta_N$ is neither real
nor of the periodic-lattice law). Clause (a) as executed reads
$\|\mathcal U_N-\mathcal U\|_{L^2(\mathbb P)}\le C(1+t)(1+L)\eps_N(1+\log(1/\eps_N))^{3/2}$ under the standing
hypothesis $tL\ge1$ (V2, Theorem 4.3). Recorded as an erratum in V4, Remark 1.1.
\end{statement}

\begin{statement}[A5: uniform-in-$N$ Nelson bounds on the space--time torus]
\label{st:A5}
\emph{[exec]} For every $p<\infty$ and $t>0$:
$\sup_N\|e^{-\mathcal U_N}\|_{L^p(\mu^{\rm per}_{t,L,N})}<\infty$, and the same for the
mean-shifted interactions of \eqref{eq:meanshift} locally uniformly in $p$-shift.
\emph{Proof pattern:} verbatim the Part I Appendix C stability argument with the
mollification scale $\kappa$ replaced by $\eps_N^{-1}$: the Wick lower bound
${:}P{:}_c\ge-b(1+c)^{p'}$ with $c=c_{N,L}\sim\tfrac1{2\pi}\log\eps_N^{-1}$, the $L^2$
ultraviolet difference between scales $N$ and $N'$ (part (a) computations), Nelson
hypercontractivity on the Gaussian space, and the double-exponential tail with
$2q=e^{L/(4p')}$. Nothing in that argument used noncompactness of space or vacuum
boundary conditions; the periodic covariance obeys the same two Bessel-type bounds.

\emph{Amended (14 Aug 2026):} as executed (V2, Theorem 7.2) the bound is stated for
$t\in[t_0,T]$ under the standing convention $t_0L\ge1$, with a constant depending also on
the coarse scale through $\eps_M$; the unrestricted range printed above is narrower in
execution. The gate consumes A5 only at $t\ge t_*=\max(1,1/L)$, so nothing is lost.
\end{statement}

\begin{statement}[A6: fixed-$t$ convergence of the thermal functionals]\label{st:A6}
\emph{[exec]} For every fixed $t\ge t_0$, $M$, and $\xi$:
$F_N(t;\xi)\to F_\infty(t;\xi)$. \emph{Proof shape (all pieces above):} write both
sides via A2 as
$\E[e^{i\Phi(J^{(\bullet)}_q)}\,e^{-\mathcal U_\bullet-\Delta\mathcal U_{\bullet,p}}]
\cdot G^{(0)}_\bullet(t;\xi)$ normalized by the sourceless expectations; the free
prefactors converge by A3; the exponents converge in every $L^p$ ($p<\infty$) by A4
(a,c) plus Gaussian hypercontractivity; uniform integrability of
$e^{-\mathcal U_N}$ by A5; conclude by Vitali (the same convergence pattern as Part I,
Proposition on path regularity, Step 2). The $q$-source $J_q$ is an honest $H^{-1}$
element (sharp-time \emph{position} sources are admissible:
$\int dk_0\,(k_0^2+\gamma^2)^{-1}<\infty$), so no obstruction arises on the position
side --- Lemma~\ref{lem:nogo} concerns momentum only.
\end{statement}

\begin{statement}[A7: partition-function convergence and the derived eventual gap]
\label{st:A7}
\emph{[exec]} For every $t>0$:
$Z_N(t)/Z^{(0)}_N(t)\to Z_\infty(t)/Z^{(0)}_\infty(t)$ (this is A6 at $\xi=0$,
sourceless), and $Z^{(0)}_N(t)\to Z^{(0)}_\infty(t)$ (explicit:
$\prod_{k\in\Gamma_N}(2\sinh(t\gamma_N(k)/2))^{-1}
\to\prod_{k\in\Gamma_\infty}(2\sinh(t\gamma(k)/2))^{-1}$, by the dispersion lemma and
the summable tail $e^{-t\gamma(k)}$).
\emph{Amended (12 Aug 2026):} the displayed free-factor limit is vacuous as
written --- both sides carry the divergent ground-energy factor and the
right-hand infinite product tends to $0$, since $\sum_k\gamma(k)=\infty$. A7 is
executed in the \emph{normal-ordered} normalization, where the free factor is
$\prod_{k\in\Gamma_N}(1-e^{-t\gamma_N(k)})^{-1}\to Z_0(t)$; see V4,
Remark 1.1(2) and Lemma 5.1. Every quantity the route consumes (eigenvalue
differences, $D_N$, gaps) is invariant under the $c$-number shift.

The conclusions must be read in that same normalization; in the \emph{un}-shifted
symbols two of them are false, since
$e_0(N)=\tfrac12\sum_{k\in\Gamma_N}\gamma_N(k)\ge r_Nm\to\infty$ forces
$Z_N(t)\to0\ne Z_\infty(t)$ and $E_j(N)\to+\infty$. Writing
$\widetilde Z_N(t):=\Trace e^{-t(K_{N,L}-e_0(N))}$ and
$\widetilde E_j(N):=E_j(N)-e_0(N)$: for every $t\ge t_*$ ($t_*>0$ fixed --- see
hypothesis (i) of Lemma~\ref{lem:laplace} and V2, the remark following Theorem 4.3, the interaction
comparison being unavailable as $t\downarrow0$) one has
$\widetilde Z_N(t)\to Z_\infty(t)=\Trace e^{-tK_L}$; with the uniform lower bound
$\inf_N\widetilde E_0(N)\ge-c$ (from A5 via $e^{-t\widetilde E_0(N)}\le
\widetilde Z_N(t)\le\widetilde Z^{(0)}_N(t)\,\|e^{-\mathcal U_N}\|$-type bounds),
Lemma~\ref{lem:laplace} applied to
$\widetilde\nu_N=\sum_j\delta_{\widetilde E_j(N)}$ yields
$\widetilde E_j(N)\to\mu_j$, the \emph{derived} eventual uniform gap
$E_1(N)-E_0(N)\to\gamma_L>0$, and $D_N(t_1)\to D_\infty(t_1)<\infty$ --- the last
two being invariant under the $c$-number shift, and being the only conclusions of
A7 that A8 consumes.
\end{statement}

\begin{statement}[A8: assembly]\label{st:A8}
\emph{[exec, short]} A6 + A7 + A1 verify hypotheses (i)--(iii) of
Lemma~\ref{lem:interchange} for $a_N=F_N$, $\ell_N=\omega_N(W_N(R_M^N\xi))$,
$\ell_\infty=\omega^{\rm ct}_{L,\lambda}(W_{\rm ct}(R_M^\infty\xi))$ (the spectral tail
bound preceding \S\ref{sec:routeA}.2 supplies \eqref{eq:tail} on both sides; the
$N=\infty$ instance of \eqref{eq:tail} uses A1's gap). Hence \eqref{eq:II11} holds
pointwise for every $M,\xi$; the proved Weyl criterion and identification corollary
then deliver the full-sequence, identified, fixed-coarse theorem. No statement beyond
A1--A7 is consumed.
\end{statement}

\subsection{What Route A does not need}
No $N$-uniform gap as an input (A7 derives an eventual one); no interacting
Lieb--Robinson estimate; no spatial cutoff (torus, $g\equiv1$); no slab limits or
boundary conditions (compact time circle); no uniformity of \eqref{eq:II11} in $\xi$
(Remark~\ref{rem:pointwise}); no comparison of the full lattice and continuum
\emph{measures} --- Statement A4's coupling compares only fourth-moment functionals,
which is essential: the per-mode Gaussian relative densities diverge near the band
edge (the thermal shadow of the finest-scale obstruction), so any route through
measure-level absolute continuity would fail.

\section{Route B (slab diagonal), recorded as fallback}\label{sec:routeB}

The previously envisaged route represents $\omega_N$ by $T\to\infty$ slabs and
compares slabs at fixed $T$. Its defect, stated precisely: the error
$|\omega_N(W)-\mathrm{slab}_N(T;W)|$ is controlled only by the \emph{lattice} gap
$\gamma_{N,L}$, so interchanging $T\to\infty$ with $N\to\infty$ requires exactly the
uniformity that Route A derives in A7 --- but without the trace structure, no
partition-function argument is available to derive it, and one is thrown back on an
$N$-uniform cluster-expansion input (the Battle--Federbush audit was negative). Route
B is therefore strictly weaker than Route A at current knowledge and is kept only as
a fallback should the trace-class package A1 fail to close (it will not: the free
trace is elementary and the interacting dressing is A5).

\section{Imports to acquire, and the numeric checkpoint}\label{sec:imports}

\subsection{Acquisition list}
For A1: H{\o}egh-Krohn, \emph{Relativistic quantum statistical mechanics in
two-dimensional space-time} (CMP 38 (1974) 195--224); Figari--H{\o}egh-Krohn--Nappi
(CMP 44 (1975) 265--278). Both to be added to \texttt{refs/part\_ii/} with page-image
dossier entries if used as imports. Everything else in A2--A8 is self-contained or
already proved in the Part II documents (free-rate lemmas) and Part I Appendix C
(stability pattern).

\subsection{Numeric falsification checkpoint (before execution)}
A \texttt{numpy} script, in the style of \texttt{m2\_ii\_checkpoint.py}:
\begin{enumerate}
\item verify the A3 closed forms against exact diagonalization of small free lattice
models (one and two modes, then $r_N\le64$), including the momentum component and the
CCR phase of A2(a);
\item measure the thermal fixed-coarse rate exponent at $t\in\{1,2,4\}$ against the
vacuum exponent $\theta$ (they should agree within the envelope constants);
\item for a two-site anharmonic toy ($\Lambda$ of size $2$, quartic coupling), verify
the interchange mechanism numerically: $E_j(N)$-convergence under mode-doubling, the
$D_N(t)$ tail, and $F_N(t)\to\omega_N$ rates --- a sanity check of
Lemmas~\ref{lem:interchange}--\ref{lem:laplace} in vivo;
\item compute $\delta_N=c_{N,L}-c^{\rm ct}_{L,N}$ for the A4b dictionary numerically
before the analytic evaluation, to fix the expected constant.
\end{enumerate}
Falsification criteria: any mismatch in (1) kills A2/A3 as formulated; a
$t$-dependent drift in (2) kills the claimed $t$-uniformity of A3(b); failure of
eigenvalue convergence in (3) would signal a flaw in the A7 logic.

\section{Execution phases and gates}\label{sec:phases}

\begin{center}
\begin{tabular}{P{1.6cm}P{6.6cm}P{5.6cm}}
\toprule
phase & content & gate \\
\midrule
V0 & numeric checkpoint (\S\ref{sec:imports}.2) & all four criteria pass; exponents
recorded \\
V1 & A2 (trace representations) + A3 (free thermal formulas and rate) & proofs
complete; A3(a) matches V0 item (1) to rounding \\
V2 & A4 (interaction $L^2$ + dictionary $\delta_\infty$; four-edge,
real-surrogate, and shifted fundamental-band repairs incorporated) + A5 (uniform Nelson) &
constants explicit; $\delta_\infty$ matches V0 item (4) \\
V3 & A1 (torus package) & full proof executed; finite-mode normalization and
periodic tails explicit; gap $\gamma_L>0$ certified \\
V4 & A6 + A7 + A8 (assembly) & (II.11) closed; Weyl criterion invoked; adversarial
pass on the assembled document \\
\bottomrule
\end{tabular}
\end{center}

Standing rules carried over from Part I: transcription discipline for every display
reused from the free-rate document and Appendix C; every import through a dossier
entry with page images; a claim ledger in the final document; an independent
adversarial pass before the gate is declared closed.

\section{Claim ledger of this scoping pass}

\begin{longtable}{P{5.2cm}P{2.2cm}P{6.4cm}}
\toprule
claim & status & ground \\
\midrule
\endhead
Pointwise-in-$\xi$ convergence suffices for the gate & proved elsewhere &
Weyl-reduction document, Thm.~1.1/Cor.~3.1; Remark~\ref{rem:pointwise} \\
Sharp-time momentum sources inadmissible & \textbf{proved here} &
Lemma~\ref{lem:nogo}, display \eqref{eq:nogo} \\
Interchange of limits under (i)--(iii) & \textbf{proved here} &
Lemma~\ref{lem:interchange} \\
Spectral transfer from partition functions & \textbf{proved here} &
Lemma~\ref{lem:laplace} \\
Free thermal formulas and thermal rate & scheduled [exec] & Statement A3; V0/V1 \\
Trace representations with marked slice & scheduled [exec] & Statement A2; V1 \\
Interaction $L^2$ comparison + renormalization dictionary & scheduled [exec] &
Statement A4; V2 \\
Uniform Nelson bounds (torus, all $N$) & scheduled [exec] & Statement A5 (Part I
App.~C pattern); V2 \\
Continuum torus package & import/exec & Statement A1; V3 \\
Fixed-$t$ convergence; partition functions; assembly & scheduled [exec] &
Statements A6--A8; V4 \\
Eventual uniform lattice gap & derived, not assumed & A7 via
Lemma~\ref{lem:laplace} \\
Measure-level absolute continuity route & rejected & band-edge divergence of
per-mode relative densities (\S\ref{sec:routeA}.3) \\
Route B (slab diagonal) & fallback only & \S\ref{sec:routeB} \\
\bottomrule
\end{longtable}

\end{document}
